\documentclass{article}
\usepackage{/globals/boxes}
\usepackage{/globals/de}
\usepackage{/globals/anki}
\ankiprefix{epks}
\ankitopic{EPKs}
\ankishow
\begin{document}
\section{Ereignisgesteuerte Prozessketten} 
\emph{Ereignisgesteuerte Prozessketten} (\emph{EPK}s) modellieren Geschäftsprozesse.
\ankinote{name}{Name}{Ereignisgesteuerte Prozessketten}

Hier werden zeitlich hintereinander die auszuführenden Funktionen eines Prozesses dargestellt, dazwischen die Ereignisse die aus den Funktionen folgen oder sie auslösen.
Funktionen und Ereignisse wechseln sich immer ab.
\ankinote{aufbau}{Aufbau}{Ereignis -> Funktion -> Ereignis}

Mit Operatoren können mehrere Verzweigungen des Prozesses parallel oder alternativ ausgeführt werden.

\subsection{Symbole}
\begin{description}
 \item[Ereignis] Ein Ereignis, dargestellt durch ein längliches Sechseck, beschreibt einen Zustand der eingetreten ist.
\ztodo{tikz}
\ankinote{symbol-ereignis}{Symbol Ereignis}{längliches Sechseck, Zustand}
 \item[Funktion] Eine Funktion wird durch ein abgerundetes Rechteck dargestellt.
\ztodo{tikz}
\ankinote{symbol-funk}{Symbol Funktion}{Abgerundetes Rechteck}
 \item[Operatoren] Operatoren wird durch einen Kreis mit den $\land$- und $\lor$-Symbolen oder dem Wort \texttt{xor} dargestellt.
\ztodo{tikz}
\ankinote{symbol-op}{Symbol Operatoren}{Symbol in Kreis}
 \item[Prozesswegweiser] Ein Prozesswegweiser verweist auf einen anderen Prozess welcher nun gestartet werden kann. Er wird durch eine Funktion vor einem leicht versetzten Ereignis dargestellt.  
\ztodo{tikz}
\ankinote{symbol-prwgws}{Symbol Prozesswegweiser}{Funktion vor Ereignis, zeigt auf Prozess}
 \item[Organisationseinheit] Eine Organisationseinheit ist ein Oval mit einer vertikalen Linie auf der linken Seite.
\ztodo{tikz}
\ankinote{symbol-org}{Symbol Organisationseinheit}{Oval mit vertikaler Linie}
 \item[Dokument, Objekt] Eine Dokument ist ein Rechteck mit einer welligen Unterseite.
\ztodo{tikz}
\ankinote{symbol-doc}{Symbol Dokument}{Rechteck, wellige Linie}
 \item[Anwendungssystem] Ein Anwendungssystem ist ein Rechteck mit zwei vertikalen Linien, eine auf beiden Seiten.
\ztodo{tikz}
\ankinote{symbol-system}{Symbol Anwendungssystem}{Rechteck, Linien}
 \item[Datenbank] Eine Datenbank wird durch eine Skizze eines Zylinders dargestellt.
\ztodo{tikz}
\ankinote{symbol-datenbank}{Symbol Datenbank}{Zylinder}
\end{description}
Pfeile bei denen Informationen fließt, \zb aus oder in ein Dokument oder einer Datenbank, zeigen in die Richtung des Informationsflusses.
\ankinote{info}{Informationsfluss}{Pfeil Richtung Infofluss}
\end{document}